\chapter{\texorpdfstring{TLA\textsuperscript{+}}{TLA+} model of the allocation CRDT}\label{sec:appendix_core_tla}

This appendix reproduces the complete contents of
\texttt{model/tlaplus/Core.tla} from the Cy reference implementation.
Core.tla is the primary source of truth for the semantics of the topic
allocation CRDT described in section~\ref{sec:session_crdt}; where this
specification and Core.tla disagree, Core.tla is authoritative.
The \verb|Check_| statements have been removed from this transcript.

\begin{minted}[fontsize=\scriptsize]{text}
------------------------------ MODULE Core ------------------------------
\* Key operators defined on Cyphal named topics. This is the core part of the formal protocol specification.
\* This is pure TLA+ without PlusCal.
\*
\* The decentralized topic subject-ID allocation protocol in Cyphal solves the problem of allocating unique subject-IDs
\* per topic name. The objective is to find a bijective configuration that is conflict-free
\* (one topic name per subject-ID) and non-divergent (one subject-ID per topic name).
\*
\* There is no central coordinator. Nodes eventually find the consensus by periodically broadcasting gossip messages.
\* Each gossip message carries a single topic->subject allocation entry (nodes apply simple rules to choose which entry
\* to gossip next). Every node only keeps the entries that it needs (publishers/subscribers that the application needs
\* to use); no node is required to keep the full set of entries. Every node must be *eventually* able to receive and
\* send gossip messages for the protocol to make progress; transient inability to do so will not break the protocol
\* but may delay the consensus (resource-limited nodes may choose to drop some gossip messages to reduce processing
\* burden; this does not violate the core assumptions of the protocol as long as such nodes are able to *eventually*
\* receive each live allocation entry).
\*
\* The allocation protocol does not operate on topic names directly, substituting them with numerical hashes instead.
\* The application layer is expected to attach names to hashes but this is of no relevance to the core protocol.
\* From the standpoint of the core allocation protocol, "topic name" and "topic hash" can be used interchangeably.
\*
\* Each allocation entry contains three fields: the topic hash, the topic eviction count, and the topic age.
\* A pair of entries can be compared to each other to determine if the pair constitutes a collision
\* (different names, same subject) or a divergence (same name, different subjects);
\* other comparison outcomes are of no interest.
\*
\* The eviction count is used as a Lamport clock and hence implementations ensure that it cannot overflow. Since
\* evictions are static on a quiescent network, even relatively narrow representations (e.g., 32 bit) may suffice.
\*
\* The subject-ID assigned to a topic is defined as some function of its hash and eviction count. One possible way
\* to define the function is, assuming 64-bit unsigned arithmetics:
\*
\*      subject_id = 8192 + ((hash + evictions**2) % 8380403)
\*
\* For the background on the subject of open addressing schemes and why the specific values were chosen this way, see:
\* - https://en.wikipedia.org/wiki/Quadratic_probing
\* - https://github.com/OpenCyphal-Garage/cy/issues/12
\*
\* The eviction count of a topic is incremented whenever the topic is involved in a collision and loses arbitration.
\* Arbitration is defined as two functions, defined here as LeftWinsCollision and LeftWinsDivergence.
\* The topic age is incremented monotonically as some kind of real or logical clock, such that each node that holds
\* this topic is expected to increment its age at approximately the same rate (the protocol does not require the
\* increment rate to match accurately, but better rate matching results in faster worst-case convergence).
\* The age is used to determine which topic wins arbitration in the event of a collision or a divergence,
\* with the core idea being that older topics win arbitration to avoid disturbances to established networks.
\*
\* The topic-subject allocation table described above is a kind of CRDT (conflict-free replicated data type).
\* A CRDT only guarantees eventual convergence while the system is quiescent, which seemingly contradicts the
\* use of the age counter that increments continuously. To work around this, the protocol uses the binary
\* logarithm of topic ages for topic comparison instead of the actual age value, which naturally slows down
\* the increment rate over time, approaching quiescence.
\*
\* When a node creates a new topic (i.e., when the local application creates a new subscription or publication),
\* its age and eviction counter are set to zero; the initial subject-ID is obtained using the function above.
\* The topic is commissioned to use immediately without the need to wait for the network to converge,
\* despite the fact that it is known that the initial allocation may be conflicting or divergent.
\* To avoid the risk of data misinterpretation should a collision occur, the transport layer carries the
\* topic hash per transfer, and the transfer is rejected should the hash be incorrect.
\* Each local topic is reconfigured as necessary should the allocation protocol require so, transparently for
\* the application.
\*
\* Upon reception of a gossip, the age of the matching local topic, if there is one, is updated as:
\* local_age = max(local_age, 2**remote_age_log2).
\*
\* Pavel Kirienko <pavel@opencyphal.org>, MIT license

EXTENDS Utils
LOCAL INSTANCE Integers
LOCAL INSTANCE TLC
LOCAL INSTANCE Sequences
LOCAL INSTANCE FiniteSets

\* Computes floor(log2(x)), with the special case floor(log2(0)) == -1
RECURSIVE Log2Floor(_)
Log2Floor(x) == IF x > 1 THEN 1 + Log2Floor(x \div 2) ELSE x - 1

\* Special case: Pow2(-1) = 0
Pow2(x) == IF x > 1 THEN 2^x ELSE x + 1

FloorToPow2(x) == Pow2(Log2Floor(x))

\***********************************************************************************************************************
\* Subject-ID mapping function. The ring size is the total number of distinct subject-IDs.
\* TODO: Switch to quadratic probing: https://github.com/OpenCyphal-Garage/cy/issues/12
SubjectIDRing == 10
SubjectID(hash, evictions) == (hash + evictions) % SubjectIDRing

\***********************************************************************************************************************
\* COLLISION is when topics with different hashes land on the same subject-ID (see the subject_id mapping).
LeftWinsCollision(a, b) == IF   Log2Floor(a.age) # Log2Floor(b.age)
                           THEN Log2Floor(a.age) > Log2Floor(b.age)
                           ELSE a.hash < b.hash

\* DIVERGENCE is when topics with the same hash are found to have different eviction counters:
\* the eviction counters have to match, so we need to decide which one is the correct one.
LeftWinsDivergence(a, b) == \/  Log2Floor(a.age) > Log2Floor(b.age)
                            \/ (Log2Floor(a.age) = Log2Floor(b.age) /\ a.evictions > b.evictions)

\***********************************************************************************************************************
\* The result is Nothing if no such topic exists.
\* Topics are stored in sets because all operations on them are ordering-invariant,
\* which is a basic prerequisite for CRDT operations.
GetByHash(hash, topics)            == Get(topics, LAMBDA x: x.hash = hash)
GetBySubjectID(subject_id, topics) == Get(topics, LAMBDA x: SubjectID(x.hash, x.evictions) = subject_id)

\***********************************************************************************************************************
\* A set of topics without the specified one. Same set if the specified topic is not a member.
RemoveTopic(hash, topics) == { t \in topics : t.hash # hash }

\* Add or replace a topic into a set of topics.
\* Use this to mutate a topic state in a set.
ReplaceTopic(new, topics) == {new} \cup RemoveTopic(new.hash, topics)

\* A set of topics extended with the specified one, and the existing topics possibly altered.
\* Uniqueness is guaranteed; if the topic is in the set already, it will be modified.
\* This can also be used to model state update of the local topic table.
RECURSIVE AllocateTopic(_, _)
AllocateTopic(t, topics) ==
    LET ts == RemoveTopic(t.hash, topics)
        x == GetBySubjectID(SubjectID(t.hash, t.evictions), ts)
        Evicted(z) == [hash |-> z.hash, evictions |-> 1 + z.evictions, age |-> z.age]
    IN   IF x = Nothing             THEN {t} \cup ts
    ELSE IF LeftWinsCollision(t, x) THEN AllocateTopic(Evicted(x), {t} \cup ts)
    ELSE                                 AllocateTopic(Evicted(t), ts)          \* Retry with evictions+1

\***********************************************************************************************************************
\* Constructs a conflict-free topic set. This is meant for constructing the initial node state.
\* Each hash can occur at most once.
RECURSIVE AllocateTopics(_, _)
AllocateTopics(new, topics) ==
    IF Cardinality(new) = 0 THEN topics
    ELSE LET t == CHOOSE x \in new : TRUE
         IN AllocateTopics(new \ {t}, AllocateTopic(t, topics))

\***********************************************************************************************************************
\* Implementation of the divergence resolution rule with CRDT age merge.
LOCAL AcceptGossip_Divergence(remote, topics) ==
    LET hash == remote.hash
        local == GetByHash(hash, topics)
        \* The serialization protocol floors the remote age to pow2 before transmission.
        new_age == Max({IF local # Nothing THEN local.age ELSE 0, FloorToPow2(remote.age)})
    IN
    IF local # Nothing THEN
        IF local.evictions # remote.evictions /\ LeftWinsDivergence(remote, local)
        THEN AllocateTopic([hash|->hash, evictions|->remote.evictions, age|->new_age], topics)
        ELSE              {[hash|->hash, evictions|-> local.evictions, age|->new_age]} \cup RemoveTopic(hash, topics)
    ELSE topics

\* Implementation of the collision resolution rule.
LOCAL AcceptGossip_Collision(remote, topics) ==
    LET local == GetBySubjectID(SubjectID(remote.hash, remote.evictions), topics)
    IN IF local # Nothing /\ LeftWinsCollision(remote, local)
    THEN AllocateTopic([hash |-> local.hash, evictions |-> local.evictions + 1, age |-> local.age], topics)
    ELSE topics

\* An updated sequence of topics based on a received gossip message.
\* The sequence may be modified if the gossiped entry is observed to diverge or collide with a local entry,
\* and the affected local entry loses arbitration to the gossiped one.
\* The naive implementation is as simple as:
\*
\*      AcceptGossip_Collision(remote, AcceptGossip_Divergence(remote, topics))
\*
\* which is correct and converges eventually; however, it is not optimal because it may cause a local entry
\* to lose a collision arbitration to a divergent remote entry, which will cause unnecessary eviction counter
\* increments and thus delay convergence. To avoid this, we ignore collisions if the remote entry is seen to be
\* divergent from a local entry and the local entry wins the divergence arbitration.
AcceptGossip(remote, topics) ==
    IF GetByHash(remote.hash, topics) # Nothing     \* If the topic exists locally
    THEN AcceptGossip_Divergence(remote, topics)    \* Divergence resolver will ensure collision-freedom as well
    ELSE AcceptGossip_Collision(remote, topics)     \* We don't know this topic, just ensure collision-freedom

\***********************************************************************************************************************
\* Divergent allocation detector operating on a function node_id -> {topic_0, topic_1, ..., topic_n}
\* A divergent allocation occurs when topic records with the same hash have distinct eviction counters.
\* If no divergences are found, the result is an empty set.
FindDivergent(node_topics) ==
    LET hashes == { t.hash : t \in UNION Range(node_topics) }
        flat == { [hash |-> t.hash, evictions |-> t.evictions] : t \in UNION Range(node_topics) }
    IN { h \in hashes : Cardinality({ s \in flat : s.hash = h }) > 1 }

\***********************************************************************************************************************
\* Allocation collision detector operating on a function node_id -> {topic_0, topic_1, ..., topic_n}
\* A collision occurs when topic records with distinct hashes have identical subject-ID.
\* If no collisions are found, the result is an empty set.
FindCollisions(node_topics) ==
    LET flat == { [hash |-> t.hash, subject_id |-> SubjectID(t.hash, t.evictions)] : t \in UNION Range(node_topics) }
        subject_ids == { t.subject_id : t \in flat }
    IN { s \in subject_ids : Cardinality({ t \in flat : t.subject_id = s }) > 1 }

\***********************************************************************************************************************
\* Whether the network has converged to a stable state.
\* Once a stable state is reached, the topics that are involved in it will no longer be altered as long as
\* no older topics are introduced to the network (such as in the event of network departitioning).
\* Appearance of young topics will not affect the existing entries by design.
Converged(node_topics) ==
    /\ FindDivergent(node_topics) = {}
    /\ FindCollisions(node_topics) = {}

========================================================================================================================
\end{minted}

\begin{minted}[fontsize=\scriptsize]{text}
------------------------------ MODULE Utils ------------------------------
\* General-purpose utilities not specific to this problem.
\* There are a few utilities here copied over from https://github.com/tlaplus/CommunityModules (MIT license)

LOCAL INSTANCE Integers
LOCAL INSTANCE Naturals
LOCAL INSTANCE TLC
LOCAL INSTANCE Sequences
LOCAL INSTANCE FiniteSets

CONSTANT Nothing

\* From https://github.com/tlaplus/CommunityModules (MIT license)
PrintVal(id, exp) == Print(<<id, exp>>, TRUE)
IsInjective(f) == \A a,b \in DOMAIN f : f[a] = f[b] => a = b
SeqToSet(s) == {s[i] : i \in DOMAIN s}
SetToSeq(S) == CHOOSE f \in [1..Cardinality(S) -> S] : IsInjective(f)
SetToSeqs(S) == LET D == 1..Cardinality(S) IN { f \in [D -> S] : \A i,j \in D : i # j => f[i] # f[j] }
Range(f) == { f[x] : x \in DOMAIN f }

Min(S) == CHOOSE s \in S : \A t \in S : s <= t
Max(S) == CHOOSE s \in S : \A t \in S : s >= t

\* The first element in a sequence where test(e) is TRUE; Nothing if the test is FALSE for all elements.
FirstMatch(haystack, test(_)) ==
    LET i == CHOOSE i \in 1..(Len(haystack)+1): (i > Len(haystack)) \/ test(haystack[i])
    IN IF i > Len(haystack) THEN Nothing ELSE haystack[i]

\* The set element where test(e) is TRUE; Nothing if the test is false for all elements.
Get(haystack, test(_)) == LET matches == { x \in haystack : test(x) }
                          IN IF matches = {} THEN Nothing ELSE CHOOSE x \in matches : TRUE

\* A sequence that contains zero needles.
SeqWithout(haystack, needle) == SelectSeq(haystack, LAMBDA x: x # needle)

\* <<1, 2, 3>> ==> <<2, 3, 1>>
SeqRotLeft(seq) == IF Len(seq) > 0 THEN Tail(seq) \o <<Head(seq)>> ELSE seq

\* Converts {<<k1, v1>>, <<k2, v2>>, ...} into a function [k |-> v]. Keys must be unique.
FunFromTupleSet(S) == [
    k \in {p[1] : p \in S} |-> CHOOSE v \in {r[2] : r \in S} : <<k, v>> \in S
]
========================================================================================================================
\end{minted}
